\documentclass[10pt,a4paper]{amsart}
\usepackage[margin=19mm]{geometry}
\usepackage{amsmath,amssymb,mathtools}
\usepackage[scr=rsfs,cal=euler]{mathalfa}
\usepackage{xfrac}
\usepackage[unicode,hidelinks,bookmarks=false]{hyperref}
\newcommand{\ie}{i.e.\ }
\newcommand{\cf}{cf.\ }
\newcommand{\resp}{respectively\ }
\newcommand{\Subsection}{\subsection}
\providecommand{\nobreakdash}{\nobreak\hbox{-}}
\setlength{\emergencystretch}{2em}
\multlinegap0pt
\usepackage{amssymb}
\usepackage{xcolor}
\newtheorem{lemma}{Lemma}
\title{On the binary digits of the Erd\H{o}s-Borwein constant}
\author{John M. Campbell}
\date{Source: arXiv:2605.24160v1 (CC BY 4.0)}
\begin{document}
\maketitle

\noindent\emph{Source and licence.} On the binary digits of the Erd\H{o}s-Borwein constant. Version arXiv:2605.24160v1 (CC BY 4.0), distributed by arXiv under the Creative Commons Attribution 4.0 International licence, http://creativecommons.org/licenses/by/4.0/, as recorded in the arXiv licence metadata for this exact version and shown on the abstract page https://arxiv.org/abs/2605.24160v1. Full licence notice: \texttt{licenses/arxiv-2605.24160-CC-BY-4.0.txt}.\par\smallskip

\noindent\textit{Editorial scope.} Counted unit: the article's lemma taillemma (source0.tex lines 320--340). The article's complete original proof of it is delivered with it (source0.tex lines 342--470, one \texttt{proof} environment of 129 lines). No mathematics is written, repaired or re-derived: the statement and the proof are the article's own lines, in the article's own notation, and no retained line is substituted or re-flowed.  Retained uncounted as the source-local context the proof needs: lines 224--238; lines 226--228; lines 230--232.  Nothing was omitted from any retained span, so the package carries no omission record.  No retained line contains a citation, so the wrapper carries no bibliography and no bibliography entry.  No retained line contains a figure environment, an image-inclusion command or drawing code, so no image is delivered and the package declares no asset dependency.  Editorial formatting only, in addition: an AMS \texttt{amsart} preamble that restates the article's own declarations and macros used by the retained lines, namely the environments \texttt{lemma} on separate counters, together with the standard packages those lines need.  The retained environments print in this document as Lemma 1 in order of appearance, whereas the article prints the same items with the numbering of its own full document.  The complete native source of the article as distributed in the arXiv e-print for this exact version is bundled unchanged as \texttt{sources/201256/source0.tex}.
\begin{lemma}\label{basiclogbound}
 Let $a, A, M \in \mathbb{N}$, let $Y \geq 3$, and suppose that 
\begin{equation}\label{aA1} 
 (a, A) = 1 
\end{equation}
 and that 
\begin{equation}\label{aM1Ay} 
 a + (M - 1) A \leq Y. 
\end{equation}
 Then $$ \sum_{m=0}^{M-1} d(a + m A) \leq 2 M \left( 1 + \frac{1}{2} \log Y \right) + 2 \sqrt{Y}. $$ Moreover, if $ \sqrt{Y} \leq M \log 
 Y$, then 
\begin{equation}\label{withsqrtMlog}
 \sum_{m=0}^{M-1} d(a + m A) \leq 5 M \log Y. 
\end{equation}
\end{lemma}

\begin{equation} 
 (a, A) = 1 
\end{equation}

\begin{equation} 
 a + (M - 1) A \leq Y. 
\end{equation}

\begin{lemma}\label{taillemma}
 There exist absolute constants $C_{\operatorname{tail}} > 0$ and $X_0 \geq 3$ such that the following holds. For $X \geq X_0$, set 
 $\Lambda = \frac{\log X}{\log 2}$ and $k = \lfloor \Lambda^{1/10} \rfloor$
 and $L = \lfloor \Lambda^2 \rfloor$, and let $A, M \in \mathbb{N}$ and 	 $r \in \mathbb{N}_{0}$ and $Y \geq 3$, 
 and assume that 
\begin{align}
 Y & \leq 2^{L}, \label{Yleq2L} \\ 
 \sqrt{Y} & \leq M \log Y, \text{and} \label{sqrtYMlogY} \\ 
 r + (L-1) + (M-1) A & \leq Y. \label{longestassumption} 
\end{align}
 Also assume that: For a prime $p$, we have that 
\begin{equation}\label{ifpmidA}
 p \mid A \Longrightarrow p > L. 
\end{equation}
 Also assume that: For each prime divisor $p \mid A$, there is an integer $ j_p \in \{ 0, 1, \ldots, k-1 \} $ such that 
\begin{equation}\label{jpcongruence} 
 r + j_p \equiv 0 \pmod p. 
\end{equation}
 For 	 $m \in [0, M)$, set $ n_{m} = r + m A $ and $ T_{m, k} = \sum_{\ell \geq k} \frac{d(n_m + \ell)}{2^\ell}$. Then $$ \#\Big\{ m \in [0, 
 M) : T_{m, k} > 2^{-k/2} \Big\} \leq C_{\operatorname{tail}} M (\log Y) 2^{-k/2}.$$
\end{lemma}

\begin{proof}
 Choose $X_0$ to be large enough so as to guarantee that: For all $X \geq X_0$, we have that $k \geq 1$ and that $L \geq 2$ and that 
\begin{equation}\label{Lhalfk} 
 \frac{L}{2} \geq k 
\end{equation} 
and that 
\begin{equation}\label{sqrtLpower} 
 \sqrt{L} + 1 \leq 2^{L/2}.
\end{equation}
 This is possible, from the given definitions for $k$ and $L$. Rewrite $\sum_{m=0}^{M-1} T_{m, k}$ so that 
\begin{align*}
 \sum_{m=0}^{M-1} T_{m, k}
 & = \sum_{m=0}^{M-1} \sum_{\ell \geq k} \frac{d(r + m A + \ell)}{2^{\ell}} \\
 & = \sum_{\ell \geq k} 2^{-\ell} \sum_{m=0}^{M-1} d(r + \ell + m A). 
\end{align*}
 Writing 
\begin{align}
 S_1 & = \sum_{\ell = k}^{L-1} 2^{-\ell} \sum_{m=0}^{M-1} d(r + \ell + m A) \ \text{and} \label{displayS1} \\ 
 S_{2} & = \sum_{\ell \geq L} 2^{-\ell} \sum_{m=0}^{M-1} d(r + \ell + m A), \label{displayS2}
\end{align}
 i.e., with $ \sum_{m=0}^{M-1} T_{m, k} = S_{1} + S_{2}$. We proceed to bound $S_{1}$. 

 Fix 
\begin{equation}\label{ellkL}
 \ell \in [k, L), 
\end{equation}
 i.e., so that $\ell$ is among the indices in the outer sum in \eqref{displayS1}. We claim that 
\begin{equation}\label{coprimerellA}
 (r + \ell, A) = 1. 
\end{equation} 
 By way of contradiction, suppose that there exists a prime $p$ such that $p \mid A$ and $p \mid (r + \ell)$. Since $p \mid A$, this 
 \emph{same} prime is such that that there exists $j_p \in [0, k)$ such that \eqref{jpcongruence} holds. From \eqref{jpcongruence} 
 together with the congruence $ r + \ell \equiv 0 \pmod p$, we have, as a consequence, that 
\begin{equation}\label{elljp0p} 
 \ell - j_p \equiv 0 \pmod p. 
\end{equation}
 Since $\ell \geq k$ and since $j_p \leq k -1$, we have that $\ell - j_p \geq 1$. Since $\ell < L$ and since $j_p \geq 0$, we have that 
 $\ell - j_p < L$. Consequently, 
 the bounds $1 \leq \ell - j_p < L$ hold. By assumption, each prime divisor $p$ of $A$ is such that $p 
 > L$, i.e., so that $1 \leq \ell - j_p < p$, but this contradicts (via \eqref{elljp0p})
 that $p$ divides $\ell - j_p$. 
 This proves the relation in \eqref{coprimerellA}. 

 Now, set $a = r + \ell$. Recall that $r \in \mathbb{N}_{0}$. Also recall that $k \geq 1$, with $\ell \geq k$. So, the bound $a \geq 1$ 
 holds. As above, we have shown that $(a, A) = 1$ holds, i.e., so that the condition displayed in \eqref{aA1} within Lemma 
 \ref{basiclogbound} holds. We again have that $\ell \leq L-1$, recalling that we are working under the assumption that $k$ is in the 
 index set indicated in \eqref{ellkL}. With regard to the condition in \eqref{aM1Ay} associated with Lemma \ref{basiclogbound}, 
 we find that 
\begin{align*}
 a + (M-1) A & = r + \ell + (M-1) A \\ 
 & \leq r + (L-1) + (M-1) A, 
\end{align*}
 so that the assumption in \eqref{longestassumption} gives us that $a + (M - 1) A \leq Y$ holds, so that all of the fundamental 
 conditions for Lemma \ref{basiclogbound} are satisfied. Moreover, from the assumption in \eqref{sqrtYMlogY}, Lemma \ref{basiclogbound} 
 gives us (via \eqref{withsqrtMlog}) that $$ \sum_{m=0}^{M-1} d(r + \ell + m A) \leq 5 M \log Y, $$ i.e., so that $ S_{1} \leq 5 M \log Y 
 \sum_{\ell = k}^{L-1} 2^{-\ell}$, which implies that 
\begin{equation}\label{S1final} 
 S_{1} \leq 10 M (\log Y) 2^{-k}. 
\end{equation}

 We proceed to bound $S_2$. In this case, we exploit the elementary bound
\begin{equation}\label{dsqrt}
 d(N) \leq 2 \sqrt{N}, 
\end{equation}
 which follows from a pairing argument given previously. Starting with the assumption in \eqref{longestassumption}, from the above 
 assumption that $L \geq 2$, we have that $ r + (M-1) A \leq Y$, so that 
\begin{equation}\label{implyboundY}
 m \in [0, M) \Longrightarrow r + m A \leq Y. 
\end{equation}
 We henceforth let $\ell \geq L$, i.e., so that $\ell$ is an index associated with the outer sum in \eqref{displayS2}. Again letting $m \in [0, 
 M)$, from \eqref{implyboundY}, we write $r + \ell + m A \leq Y + \ell$, so that the divisor bound in \eqref{dsqrt} gives us that $ d(r + 
 \ell + m A) \leq 2 \sqrt{Y + \ell}$, i.e., so that 
\begin{equation}\label{S2elementary}
 S_{2} \leq 2 M \sum_{\ell \geq L} 2^{-\ell} \sqrt{Y + \ell}. 
\end{equation}
 Using the inequality $\sqrt{Y + \ell} \leq \sqrt{Y} + \sqrt{\ell}$, we have that 
\begin{equation}\label{detailsforS2}
 \sum_{\ell \geq L} 2^{-\ell} \sqrt{Y + \ell}
 \leq \sqrt{Y} \sum_{\ell \geq L} 2^{-\ell} + \sum_{\ell \geq L} 2^{-\ell} \sqrt{\ell}, 
\end{equation}
 by evaluating the first series on the right of \eqref{detailsforS2} in closed form, and by using the bound on $Y$ in \eqref{Yleq2L}, we 
 find that 
\begin{equation}\label{toreindex}
 \sum_{\ell \geq L} 2^{-\ell} \sqrt{Y + \ell}
 \leq 2^{1 - L/2} + \sum_{\ell \geq L} 2^{-\ell} \sqrt{\ell}. 
\end{equation}
 We then apply a reindexing argument to the right-hand series in \eqref{toreindex}, writing
\begin{equation}\label{secondsqrtineq}
 \sum_{\ell \geq L} 2^{-\ell} \sqrt{\ell} = 
 2^{-L} \sum_{t \geq 0} 2^{-t} \sqrt{L + t}, 
\end{equation}
 and we again exploit an elementary inequality of the form $\sqrt{L + t} \leq \sqrt{L} + \sqrt{t}$, with \eqref{secondsqrtineq} giving us 
 that $ \sum_{t \geq 0} 2^{-t} \sqrt{L + t} \leq \sqrt{L} \sum_{t \geq 0} 2^{-t} + \sum_{t \geq 0} 2^{-t} \sqrt{t} 
 \leq C_0 \big( \sqrt{L} + 1 \big)$, 
 for an absolute constant $C_0$. 
 Exploiting the bound 
$\sqrt{L} + 1 \leq 2^{L/2}$ in \eqref{sqrtLpower}, we obtain that 
\begin{equation}\label{aftersqrtLpow}
 \sum_{\ell \geq L} 2^{-\ell} \sqrt{\ell} 
 \leq C_0 2^{-L} \big( \sqrt{L} + 1 \big) \leq C_{0} 2^{-L/2}. 
\end{equation}
 From \eqref{toreindex} and \eqref{aftersqrtLpow}, we find that 
 $ \sum_{\ell \geq L} 2^{-\ell} \sqrt{Y + \ell} \leq 2^{1 - L/2} + C_{0} 2^{-L/2}$, i.e., so that $ \sum_{\ell \geq L} 2^{-\ell} \sqrt{Y + \ell} 
 \leq C_{1} 2^{-L/2} $ for an absolute constant $C_{1}$. Exploiting the bound in \eqref{Lhalfk}, we obtain that $ \sum_{\ell \geq L} 
 2^{-\ell} \sqrt{Y + \ell} \leq C_{1} 2^{-k}$, i.e., so that 
\begin{equation}\label{S2bound} 
 S_{2} \leq 2 C_{1} M 2^{-k}. 
\end{equation}
 Since $Y \geq 3$, we obtain from \eqref{S2bound} that 
\begin{equation}\label{S2final} 
 S_2 \leq 2 C_1 M (\log Y) 2^{-k}. 
\end{equation}
 From the bounds for $S_1$ and $S_2$ in \eqref{S1final} and \eqref{S2final}, we find that there is an absolute constant $C_{2}$ such that 
\begin{equation}\label{finalsumT} 
 \sum_{m=0}^{M-1} T_{m, k} \leq C_{2} M (\log Y) 2^{-k}. 
\end{equation}
 We proceed to set
 $ \mathcal{B} = \big\{ 0 \leq m < M : T_{m, k} > 2^{-k/2} \big\}$. 
 For $m \in \mathcal{B}$, we have that $T_{m, k} > 2^{-k/2}$, so that 
\begin{equation}\label{implicitMarkov} 
 \sum_{m \in \mathcal{B}} T_{m, k} \geq \sum_{m \in \mathcal{B}} 2^{-k/2} = | \mathcal{B} | 2^{-k/2}. 
\end{equation}
 From \eqref{implicitMarkov}, we proceed to write
\begin{equation}\label{afterMarkov}
 |\mathcal{B}| 2^{-k/2} \leq \sum_{m \in \mathcal{B}} T_{m, k} \leq \sum_{m=0}^{M-1} T_{m, k}, 
\end{equation}
 so that \eqref{finalsumT} and \eqref{afterMarkov} together give us the desired bound, 
 writing $C_{\operatorname{tail}} = C_{2}$. 
\end{proof}

\end{document}
